\documentclass[10pt,a4paper]{amsart}
\usepackage{amsmath,amssymb,amsthm,graphicx}
% Retained preamble for the Stacks Project chapter morphisms.tex (Morphisms of Schemes).
% The chapter file pinned-original-source.tex carries only \input{preamble} on its first line, so no
% preamble text lives inside the chapter itself.  The Stacks Project uses one single shared preamble for
% every chapter; the snapshot stages it verbatim in sources/stacks-chow/preamble.tex.  The block below is
% copied mechanically, line for line, from that file's lines 14-18 (the Xy-pic setup) and lines 154-217
% (the native theorem declarations together with the mathematical macros).  Nothing was added, removed or
% redefined.  The wrapper class and the amsmath/amssymb/amsthm/graphicx base packages are supplied by the
% assembler, so the class/\externaldocument part of the shared preamble is not repeated here.
% For commutative diagrams we use Xy-pic
\usepackage[all]{xy}

% We use 2cell for 2-commutative diagrams.
\xyoption{2cell}
% Theorem environments.
%
\theoremstyle{plain}
\newtheorem{theorem}[subsection]{Theorem}
\newtheorem{proposition}[subsection]{Proposition}
\newtheorem{lemma}[subsection]{Lemma}

\theoremstyle{definition}
\newtheorem{definition}[subsection]{Definition}
\newtheorem{example}[subsection]{Example}
\newtheorem{exercise}[subsection]{Exercise}
\newtheorem{situation}[subsection]{Situation}

\theoremstyle{remark}
\newtheorem{remark}[subsection]{Remark}
\newtheorem{remarks}[subsection]{Remarks}

\numberwithin{equation}{subsection}

% Macros
%
\def\lim{\mathop{\mathrm{lim}}\nolimits}
\def\colim{\mathop{\mathrm{colim}}\nolimits}
\def\Spec{\mathop{\mathrm{Spec}}}
\def\Hom{\mathop{\mathrm{Hom}}\nolimits}
\def\Ext{\mathop{\mathrm{Ext}}\nolimits}
\def\SheafHom{\mathop{\mathcal{H}\!\mathit{om}}\nolimits}
\def\SheafExt{\mathop{\mathcal{E}\!\mathit{xt}}\nolimits}
\def\Sch{\mathit{Sch}}
\def\Mor{\mathop{\mathrm{Mor}}\nolimits}
\def\Ob{\mathop{\mathrm{Ob}}\nolimits}
\def\Sh{\mathop{\mathit{Sh}}\nolimits}
\def\NL{\mathop{N\!L}\nolimits}
\def\CH{\mathop{\mathrm{CH}}\nolimits}
\def\proetale{{pro\text{-}\acute{e}tale}}
\def\etale{{\acute{e}tale}}
\def\QCoh{\mathit{QCoh}}
\def\Ker{\mathop{\mathrm{Ker}}}
\def\Im{\mathop{\mathrm{Im}}}
\def\Coker{\mathop{\mathrm{Coker}}}
\def\Coim{\mathop{\mathrm{Coim}}}

% Boxtimes
%
\DeclareMathSymbol{\boxtimes}{\mathbin}{AMSa}{"02}

%
% Macros for moduli stacks/spaces
%
\def\QCohstack{\mathcal{QC}\!\mathit{oh}}
\def\Cohstack{\mathcal{C}\!\mathit{oh}}
\def\Spacesstack{\mathcal{S}\!\mathit{paces}}
\def\Quotfunctor{\mathrm{Quot}}
\def\Hilbfunctor{\mathrm{Hilb}}
\def\Curvesstack{\mathcal{C}\!\mathit{urves}}
\def\Polarizedstack{\mathcal{P}\!\mathit{olarized}}
\def\Complexesstack{\mathcal{C}\!\mathit{omplexes}}
% \Pic is the operator that assigns to X its picard group, usage \Pic(X)
% \Picardstack_{X/B} denotes the Picard stack of X over B
% \Picardfunctor_{X/B} denotes the Picard functor of X over B
\def\Pic{\mathop{\mathrm{Pic}}\nolimits}
\def\Picardstack{\mathcal{P}\!\mathit{ic}}
\def\Picardfunctor{\mathrm{Pic}}
\def\Deformationcategory{\mathcal{D}\!\mathit{ef}}

% --- appended 2026-10-12 (Chief): the project own definitions of the reference/slogan/history
% display environments, copied mechanically from the shared Stacks preamble lines 8-12
% (\usepackage{verbatim} plus the three \newenvironment lines, which the chapter preamble keeps
% verbatim from the project preamble). Nothing added, removed or redefined: these environments
% carry no mathematical content and the project itself typesets them as comments. Without these
% three lines a counted statement whose own text uses them fails with Environment slogan undefined.
\usepackage{verbatim}
\newenvironment{reference}{\comment}{\endcomment}
\newenvironment{slogan}{\comment}{\endcomment}
\newenvironment{history}{\comment}{\endcomment}
\usepackage{cleveref}
\crefname{originaldefinition}{original definition}{original definitions}
\Crefname{originaldefinition}{Original definition}{Original definitions}
\makeatletter
\newlabel{algebra-definition-quasi-finite}{{01TD}{source}{}{sourcepointer-1}{}}
\newlabel{algebra-definition-quasi-finite@cref}{{[originaldefinition][0][]01TD}{[1][0][]source}{}{}{}}
\makeatother
\crefname{originallemma}{original lemma}{original lemmas}
\Crefname{originallemma}{Original lemma}{Original lemmas}
\makeatletter
\newlabel{lemma-base-change-finite-type}{{01T4}{source}{}{sourcepointer-2}{}}
\newlabel{lemma-base-change-finite-type@cref}{{[originallemma][0][]01T4}{[1][0][]source}{}{}{}}
\makeatother
\crefname{originallemma}{original lemma}{original lemmas}
\Crefname{originallemma}{Original lemma}{Original lemmas}
\makeatletter
\newlabel{lemma-locally-finite-type-characterize}{{01T2}{source}{}{sourcepointer-3}{}}
\newlabel{lemma-locally-finite-type-characterize@cref}{{[originallemma][0][]01T2}{[1][0][]source}{}{}{}}
\makeatother
\crefname{originallemma}{original lemma}{original lemmas}
\Crefname{originallemma}{Original lemma}{Original lemmas}
\makeatletter
\newlabel{lemma-ubiquity-Jacobson-schemes}{{02J6}{source}{}{sourcepointer-4}{}}
\newlabel{lemma-ubiquity-Jacobson-schemes@cref}{{[originallemma][0][]02J6}{[1][0][]source}{}{}{}}
\makeatother
\title{Four equivalent descriptions of quasi-finiteness at a point in terms of the fibre}
\author{The Stacks Project Authors}
\date{stacks-project:morphisms.tex:local-snapshot-20261012; GNU FDL 1.2 or later}
\begin{document}
\maketitle
\noindent\textit{Editorial note.} All delivered mathematical text is the author's own, copied from the pinned Stacks Project chapter file pinned-original-source.tex; no statement or proof step was rewritten, reordered or removed.\par\medskip
\noindent\textit{Editorial note.} Following the multi-claim convention for this source, source\_url is the source-level URL https://stacks.math.columbia.edu/ and the source label lemma-quasi-finite-at-point-characterize is carried in claim\_key and here rather than in a per-claim tag URL.\par\medskip
\noindent\textit{Editorial note.} The preamble is the chapter's shared Stacks Project preamble, retained as sources/stacks-morphisms/retained-preamble.tex: a line-for-line excerpt of the shared Stacks preamble staged with this snapshot. No macro was invented or redefined, and this card adds no preamble text of its own.\par\medskip
\noindent\textit{Editorial note.} Statement span 3663-3681 and proof span 3683-3727 come from the halfA pool card operations/pool-cards/stacks-morphisms-halfA.json and were re-read in the pinned chapter file before this card was written.\par\medskip
\noindent\textit{Editorial note.} Re-checked by hand before carding: the lane's thin-proof prefilter flagged this candidate, and reading the original shows the hit stems only from short incidental phrases such as 'It is clear that' in the closing paragraph, while the body is a full equivalence chain (isolation, Jacobson property, fibre ring dimension, affine charts). Per the fleet rule of 2026-10-14 the hit on a proof of more than 800 bytes was re-read by hand and the reason is recorded here.\par\medskip
\noindent\textit{Editorial note.} The proof refers to the algebra definition of quasi-finiteness at a prime and to the characterization of locally finite type morphisms; both references are part of the counted span and are handled as navigation pointers.\par\medskip
\noindent\textit{Editorial note.} 4 cross-references occur inside the counted spans; they are listed in external\_references as navigation pointers only.\par\medskip
\noindent\textit{Allowed standard prerequisites (editorial).}\begin{itemize}
\item The definition of quasi-finiteness at a point by the existence of affine neighbourhoods in which the point is isolated in the fibre. Used at proof lines 3683-3687 and at the end of the proof. Supplies the first step and the final step of the equivalence chain.
\item Fibres of locally finite type morphisms are Jacobson schemes, cited from the same chapter with the base change statement for finite type morphisms. Cited at proof lines 3690-3693. Shows that an isolated point of the fibre contains a closed point, hence is itself closed.
\item The characterization of locally finite type morphisms by finite type ring maps on affine charts, and the description of the fibre ring as the tensor product with the residue field, together with the algebra definition of quasi-finiteness at a prime. Cited at proof lines 3701-3704, 3706-3708 and 3716-3719. Converts the closedness and non-specialization hypothesis into the vanishing of the dimension of the fibre ring at the point, which is the algebraic quasi-finiteness condition.
\end{itemize}\medskip
\ifdefined\hypertarget\hypertarget{sourcepointer-1}{}\fi
\noindent\textit{Source reference.} Original definition 01TD: focused-seeds/AGP/sources/stacks-snapshot-20261011/tags.txt, source lines 2370--2370. Same-chapter reference used in the retained author text at source line(s) 3535, 3660, 3717, 3836. It is supplied only as a navigation pointer so that the retained text compiles; the counted item delivers the author's own statement and the author's complete proof, and no part of the referenced item is copied in or substituted for the counted argument.\par\medskip
\ifdefined\hypertarget\hypertarget{sourcepointer-2}{}\fi
\noindent\textit{Source reference.} Original lemma 01T4: focused-seeds/AGP/sources/stacks-snapshot-20261011/tags.txt, source lines 2361--2361. Same-chapter reference used in the retained author text at source line(s) 2956, 3695, 3779, 3950, 4313, 6088, 9750, 9928. It is supplied only as a navigation pointer so that the retained text compiles; the counted item delivers the author's own statement and the author's complete proof, and no part of the referenced item is copied in or substituted for the counted argument.\par\medskip
\ifdefined\hypertarget\hypertarget{sourcepointer-3}{}\fi
\noindent\textit{Source reference.} Original lemma 01T2: focused-seeds/AGP/sources/stacks-snapshot-20261011/tags.txt, source lines 2359--2359. Same-chapter reference used in the retained author text at source line(s) 2717, 2739, 2801, 2863, 3019, 3035, 3157, 3348, 3613, 3707, 3868, 5197, 5282, 6702, 9325, 9390, 9440, 11421, 13422, 13686. It is supplied only as a navigation pointer so that the retained text compiles; the counted item delivers the author's own statement and the author's complete proof, and no part of the referenced item is copied in or substituted for the counted argument.\par\medskip
\ifdefined\hypertarget\hypertarget{sourcepointer-4}{}\fi
\noindent\textit{Source reference.} Original lemma 02J6: focused-seeds/AGP/sources/stacks-snapshot-20261011/tags.txt, source lines 3325--3325. Same-chapter reference used in the retained author text at source line(s) 3693, 3744, 13786. It is supplied only as a navigation pointer so that the retained text compiles; the counted item delivers the author's own statement and the author's complete proof, and no part of the referenced item is copied in or substituted for the counted argument.\par\medskip
\begin{lemma}
\label{lemma-quasi-finite-at-point-characterize}
Let $f : X \to S$ be a morphism of schemes.
Let $x \in X$ be a point. Set $s = f(x)$.
Let $X_s$ be the fibre of $f$ at $s$.
Assume $f$ is locally of finite type.
The following are equivalent:
\begin{enumerate}
\item The morphism $f$ is quasi-finite at $x$.
\item The point $x$ is isolated in $X_s$.
\item The point $x$ is closed in $X_s$
and there is no point $x' \in X_s$, $x' \not = x$
which specializes to $x$.
\item For any pair of affine opens
$\Spec(A) = U \subset X$, $\Spec(R) = V \subset S$ with
$f(U) \subset V$ and $x \in U$ corresponding to $\mathfrak q \subset A$
the ring map $R \to A$ is quasi-finite at $\mathfrak q$.
\end{enumerate}
\end{lemma}

\begin{proof}
Assume $f$ is quasi-finite at $x$. By assumption there exist opens
$U \subset X$, $V \subset S$ such that $f(U) \subset V$, $x \in U$
and $x$ an isolated point of $U_s$. Hence $\{x\} \subset U_s$ is an open
subset. Since $U_s = U \cap X_s \subset X_s$ is also open we conclude
that $\{x\} \subset X_s$ is an open subset also. Thus we conclude
that $x$ is an isolated point of $X_s$.

\medskip\noindent
Note that $X_s$ is a Jacobson scheme by
Lemma \ref{lemma-ubiquity-Jacobson-schemes}
(and
Lemma \ref{lemma-base-change-finite-type}).
If $x$ is isolated in $X_s$, i.e., $\{x\} \subset X_s$ is open,
then $\{x\}$ contains a closed point (by the Jacobson property), hence
$x$ is closed in $X_s$. It is clear that there is no point $x' \in X_s$,
distinct from $x$, specializing to $x$.

\medskip\noindent
Assume that $x$ is closed in $X_s$ and that there is no point $x' \in X_s$,
distinct from $x$, specializing to $x$. Consider a pair of affine opens
$\Spec(A) = U \subset X$, $\Spec(R) = V \subset S$ with
$f(U) \subset V$ and $x \in U$. Let $\mathfrak q \subset A$ correspond to
$x$ and $\mathfrak p \subset R$ correspond to $s$.
By Lemma \ref{lemma-locally-finite-type-characterize} the ring map
$R \to A$ is of finite type. Consider the fibre ring
$\overline{A} = A \otimes_R \kappa(\mathfrak p)$.
Let $\overline{\mathfrak q}$ be the prime of $\overline{A}$ corresponding
to $\mathfrak q$. Since $\Spec(\overline{A})$ is an open subscheme of
the fibre $X_s$ we see that $\overline{q}$ is a maximal ideal
of $\overline{A}$ and that there is no point of $\Spec(\overline{A})$
specializing to $\overline{\mathfrak q}$.
This implies that $\dim(\overline{A}_{\overline{q}}) = 0$.
Hence by
Algebra, Definition \ref{algebra-definition-quasi-finite}
we see that $R \to A$ is quasi-finite at $\mathfrak q$, i.e.,
$X \to S$ is quasi-finite at $x$ by definition.

\medskip\noindent
At this point we have shown conditions (1) -- (3) are all equivalent.
It is clear that (4) implies (1). And it is also clear that
(2) implies (4) since if $x$ is an isolated point of $X_s$
then it is also an isolated point of $U_s$ for any open $U$
which contains it.
\end{proof}

\end{document}
